\begin{filecontents*}{references.bib}
@article{rmus2025generating,
  title={Generating Computational Cognitive Models using Large Language Models},
  author={Rmus, Milena and Mathony, Marvin and Jagadish, Akshay K. and Ludwig, Tobias and Schulz, Eric},
  journal={arXiv preprint arXiv:2502.00879},
  year={2025}
}
@article{castro2025discovering,
  title={Discovering symbolic cognitive models from human and animal behavior},
  author={Castro, Pablo Samuel and Tomasev, Nenad and Anand, Ankit and Sharma, Navodita and others},
  journal={bioRxiv},
  year={2025},
  publisher={Cold Spring Harbor Laboratory}
}
@article{lueckmann2026discovering,
  title={Discovering mechanistic models of neural activity: System identification in an in silico zebrafish},
  author={Lueckmann, Jan-Matthis and Jain, Viren and Januszewski, Micha{\l}},
  journal={arXiv preprint arXiv:2602.04492},
  year={2026}
}
@article{tilbury2025ai,
  title={AI-discovered tuning laws explain neuronal population code geometry},
  author={Tilbury, Reilly and Kwon, Dabin and Haydaro{\u{g}}lu, Ali and Ratliff, Jacob and others},
  journal={bioRxiv},
  year={2025},
  publisher={Cold Spring Harbor Laboratory}
}
@article{luo2025large,
  title={Large language models surpass human experts in predicting neuroscience results},
  author={Luo, Xiaoliang and Rechardt, Akilles and Sun, Guangzhi and Nejad, Kevin K. and others},
  journal={Nature Human Behaviour},
  year={2025},
  publisher={Nature Publishing Group}
}
@article{luppi2025cognitive,
  title={Cognitive cartography of mammalian brains using meta-analysis of AI experts},
  author={Luppi, Andrea I. and Ali, Hana and Liu, Zhen-Qi and Milisav, Filip and others},
  journal={bioRxiv},
  year={2025},
  publisher={Cold Spring Harbor Laboratory}
}
@article{ghayem2026neurocontext,
  title={NeuroConText: Contrastive learning for neuroscience meta-analysis with rich text representation},
  author={Ghayem, Fateme and Meudec, Rapha{\"e}l and Dock{\`e}s, J{\'e}r{\^o}me and Thirion, Bertrand and Wassermann, Demian},
  journal={Imaging Neuroscience},
  year={2026},
  publisher={MIT Press}
}
@article{wang2025foundation,
  title={Foundation model of neural activity predicts response to new stimulus types},
  author={Wang, Eric Y. and Fahey, Paul G. and Ding, Zhuokun and Papadopoulos, Stelios and others},
  journal={Nature},
  year={2025},
  publisher={Nature Publishing Group}
}
@article{ryoo2025generalizable,
  title={Generalizable, real-time neural decoding with hybrid state-space models},
  author={Ryoo, Avery Hee-Woon and Krishna, Nanda H. and Mao, Ximeng and Azabou, Mehdi and others},
  journal={arXiv preprint arXiv:2506.05320},
  year={2025}
}
@article{wairagkar2025instantaneous,
  title={An instantaneous voice-synthesis neuroprosthesis},
  author={Wairagkar, Maitreyee and Card, Nicholas S. and Singer-Clark, Tyler and Hou, Xianda and others},
  journal={Nature},
  year={2025},
  publisher={Nature Publishing Group}
}
@article{card2025longterm,
  title={Long-term independent use of an intracortical brain-computer interface for speech and cursor control},
  author={Card, Nicholas S. and Singer-Clark, Tyler and Peracha, Hamza and Iacobacci, Carrina and Hou, Xianda and others},
  journal={bioRxiv},
  year={2025},
  doi={10.1101/2025.06.26.661591}
}
@article{collins2012cognitive,
  title={Cognitive control over learning: Creating, clustering, and generalizing task-set structure},
  author={Collins, Anne GE and Frank, Michael J},
  journal={Psychological review},
  volume={119},
  number={4},
  pages={890},
  year={2012},
  publisher={American Psychological Association}
}
@article{meta2024llama3,
  title={The Llama 3 Herd of Models},
  author={Dubey, Abhimanyu and Jauhri, Abhinav and Pandey, Abhinav and Kadian, Abhishek and others},
  journal={arXiv preprint arXiv:2407.21783},
  year={2024}
}
@article{guo2025deepseek,
  title={DeepSeek-R1: Incentivizing Reasoning Capability in LLMs via Reinforcement Learning},
  author={Guo, Daya and Yang, Dejian and Zhang, Haowei and Song, Junxiao and Zhang, Ruoyu and Xu, Runxin and others},
  journal={arXiv preprint arXiv:2501.12948},
  year={2025}
}
@article{yang2024qwen25,
  title={Qwen2. 5 technical report},
  author={Yang, An and Yang, Baosong and Hui, Binyuan and Zheng, Bo and Yu, Bowen and others},
  journal={arXiv preprint arXiv:2412.15115},
  year={2024}
}

@article{collins2017workingmemory,
  title={Working memory load strengthens reward prediction errors},
  author={Collins, Anne G. E. and Ciullo, Brittany and Frank, Michael J. and Badre, David},
  journal={The Journal of Neuroscience},
  volume={37},
  number={16},
  pages={4332--4342},
  year={2017},
  doi={10.1523/JNEUROSCI.2700-16.2017}
}
@book{kaplan1964conduct,
  title={The Conduct of Inquiry: Methodology for Behavioral Science},
  author={Kaplan, Abraham},
  year={1964},
  publisher={Chandler}
}
\end{filecontents*}

\documentclass[12pt, a4paper]{article}

\usepackage[margin=2.5cm]{geometry}
\usepackage{hyperref}
\usepackage{booktabs}
\usepackage{array}
\usepackage{parskip}
\usepackage{titlesec}
\usepackage{xcolor}
\usepackage{graphicx}
\usepackage[T1]{fontenc}
\usepackage[utf8]{inputenc}
\usepackage[round]{natbib} % Added for proper citation formatting

\hypersetup{
    colorlinks=true,
    linkcolor=blue!60!black,
    urlcolor=blue!60!black,
    citecolor=blue!60!black
}

\titleformat{\section}{\large\bfseries}{}{0em}{}[\titlerule]
\titleformat{\subsection}{\normalsize\bfseries}{}{0em}{}

\title{\textbf{Seminar Thesis on \textit{Generating Computational Cognitive Models
using Large Language Models}}\\
\large Submitted by Christian Block and YingZhu Chen}

\author{}
\date{}

\begin{document}

\maketitle
\tableofcontents

% ─────────────────────────────────────────────────────────────
\section{Executive Summary}

The pipeline proposed in \textit{Generating Computational Cognitive Models
using Large Language Models} uses Large Language Models (LLMs) to automatically generate, fit, and iteratively refine computational cognitive models as Python functions.\footnote{In the remainder of this report, we also refer the proposed pipeline as GeCCo.} Given a task description, participant behavioural data, and a template function, the pipeline prompts an LLM to propose candidate models, fits those proposals to held-out data, and refines them based on feedback derived from predictive performance.

The approach was evaluated across four cognitive domains---decision making, learning, planning, and working memory---using three open-source LLMs that differed in model family, size, and training procedure. The strongest LLM-generated models achieved lower mean BIC scores than the hand-crafted baselines in all four domains. However, the evidence was not equally strong in every case. The BIC difference was statistically significant in decision making and working memory, whereas the advantage was smaller in learning and not statistically significant in planning. The overall result is therefore best described as competitive and sometimes superior, rather than uniformly better than existing cognitive models.

% ─────────────────────────────────────────────────────────────
\section{Motivation}

Computational cognitive models provide a formal framework for understanding human behaviour. In practice, researchers hand-craft these models, fit them to behavioural data, and refine them over several cycles. This process is slow and requires expertise in cognitive science, statistical modelling, and programming. It may also be influenced by a researcher's previous exposure to particular theories and model families. As a result, researchers may focus on familiar explanations and overlook simpler or less conventional alternatives. This resembles the \emph{streetlight effect}: investigation is drawn towards areas where existing tools and theories make searching easiest. These limitations motivate an automated approach that can explore a broader range of candidate models.\\

Three LLM capabilities make them suitable for proposing cognitive models from experimental data.\footnote{The models used were Llama-3.1-Instruct-70B, DeepSeek-R1-Distill-Llama-3.1-70B, and Qwen2.5-72B-Instruct.} First, LLMs can process natural-language task descriptions and examples of behavioural data, allowing the same general interface to be used across experimental paradigms. Second, they can use the task structure and observed behaviour to propose possible psychological mechanisms. Third, they can translate these proposals into executable Python functions. Similar abilities have been used for automated statistical modelling and probabilistic program generation; \citet{rmus2025generating} apply this idea to cognitive science.

% ─────────────────────────────────────────────────────────────
\section{The GeCCo Pipeline}
Each pipeline run consists of ten sampling iterations, and the procedure is repeated across five independent runs for each cognitive domain. The process begins with a fixed prompt containing a natural-language task description, a subset of behavioural data, detailed instructions, and a domain-specific code template. From the second iteration onwards, the prompt also contains feedback from the previous round. The LLM generates three candidate cognitive models with different parameter combinations. A complete run took no more than eight hours on four Nvidia A100 GPUs, which is relatively efficient compared with search-based or evolutionary approaches that may require several days.\\

The proposed models are parsed from the response and converted into executable Python functions. Their parameters are then fitted to a held-out validation dataset that was excluded from the prompt. To reduce the risk of the optimiser becoming trapped in a local minimum, fitting is repeated from 20 random starting points.\\
The pipeline therefore alternates between LLM-based proposal generation and
external statistical evaluation: prompt $\rightarrow$ candidate Python models
$\rightarrow$ parameter fitting $\rightarrow$ BIC feedback $\rightarrow$ revised
proposals.
After fitting, the models are evaluated using the Bayesian Information Criterion (BIC), and the best-performing candidate is selected. Its code and a list of previously used parameter combinations are included as feedback in the next prompt. This guides later proposals and discourages the LLM from repeating earlier models.

After all ten iterations, the best model from the complete run is evaluated in two ways:

\begin{description}
    \item[Baseline model comparison:] The BIC of the LLM-generated model is compared with that of an established cognitive model. This assesses whether the generated model provides a better balance between fit and complexity.
    \item[Posterior predictive checks:] The LLM-generated model is used to simulate behaviour, and the simulated patterns are compared with the observed human data. These checks show whether the model reproduces important behavioural patterns, although they do not establish that its internal mechanism is correct.
\end{description}
\section{Experiments}

Two complementary metrics are used in the evaluation. BIC balances goodness of fit against model complexity by penalising the number of free parameters. When two models fit the data equally well, BIC favours the simpler one. Exceedance Probability (EXP) estimates the probability that a model is the most frequent explanation of behaviour within the tested population. Importantly, these metrics are computed offline on held-out data using standard optimisation tools, not by the LLM itself. Model selection is therefore determined by statistical evaluation rather than by the LLM's opinion of its own output.

\subsection{Experiment 1: Decision Making}

Participants chose between two products using ratings provided by four experts whose reliability differed. The baseline was the probabilistic Weighted Additive Decision model, which uses four parameters to weight the four features. Models generated by R1 and Llama achieved lower BIC scores. The best Llama model used a single \emph{discount factor} to move smoothly between three established strategies: \emph{Take the Best}, \emph{Equal Weighting}, and \emph{Weighted Additive}. Its mean BIC was 39.40, compared with 51.97 for the four-parameter baseline, and the difference was statistically significant ($p<.001$). The discount factor resembles a prior-strength parameter in Bayesian accounts of heuristics, even though no Bayesian framework was provided in the prompt. This example gives strong evidence that GeCCo can produce a model that is both simpler and better fitting than the selected baseline.

\subsection{Experiment 2: Learning}

Participants completed a two-armed bandit task in which they repeatedly chose between two options. In the full-feedback condition, they observed both the received reward and the reward they would have obtained from the unchosen option. The baseline was a four-learning-rate version of the Rescorla--Wagner model with an additional perseveration parameter for choice repetition. The best Llama-generated model obtained a lower mean BIC than the baseline (77.97 versus 85.24) and a higher EXP (0.97 versus 0.03). It introduced a separate meta-value function for forgone outcomes and replaced the perseveration parameter with a \emph{forgetting parameter} that reduced the value of unchosen options over time. This provides an alternative interpretation of repeated choice based on memory and value decay. However, the posterior predictive checks showed that the generated model matched human behaviour about as well as the baseline, so its theoretical interpretation remains a hypothesis rather than a confirmed mechanism.

\subsection{Experiment 3: Planning}

The planning experiment used a two-stage task with common and rare transitions between states and changing reward probabilities. It tests whether participants repeat previously rewarded actions or use knowledge of the transition structure to plan ahead. The baseline Hybrid model combines model-free and model-based learning using a free weighting parameter. The best R1-generated model had a slightly lower mean BIC than the Hybrid model (432.47 versus 437.38), but the difference was not statistically significant ($p=.27$). Its EXP was higher (0.85 versus 0.15), and its simulations reproduced the main human choice pattern. The generated model used two transition-dependent learning rates and omitted temporal discounting, offering a simpler account that does not require a strict division between model-free and model-based systems. Nevertheless, the non-significant BIC difference means that the evidence supports it as a plausible alternative, not as a clearly superior explanation.

\subsection{Experiment 4: Working Memory}

Participants learned by trial and error which response was correct for each visual stimulus. Blocks contained either three stimulus--action associations (low load) or six associations (high load). The baseline RL--WM model combines slow reinforcement learning with a fast but capacity-limited working-memory module. The best Llama-generated model achieved a significantly lower mean BIC than the baseline (267.25 versus 275.20, $p<.01$) and a higher EXP (0.99 versus 0.01). Instead of changing the contribution of a separate working-memory system, it assigned different reinforcement-learning rates to the low- and high-load conditions. This suggests that cognitive load may affect the learning signal itself, an interpretation that is consistent with evidence that working-memory load changes striatal reward-prediction-error signals \citep{collins2017workingmemory}. The result is theoretically interesting, but the behavioural fit alone cannot determine whether the effect is truly located in reinforcement learning rather than working memory.

\section{Control Experiments}

A mixed-effects regression predicted BIC from \emph{reasoning ability}, \emph{model family}, and \emph{model size}. Llama performed better than Qwen despite Qwen having more parameters, suggesting that model family and training procedure may matter more than size alone. The estimated benefit of using a reasoning model was small and did not reach conventional statistical significance.

The authors also removed prompt components and analysed the resulting BIC scores. All four components contributed to performance, but their estimated effects differed. Iterative feedback had the largest effect, followed by participant data, the task description, and the code template. This result supports the importance of the refinement loop rather than the initial prompt alone, although the modest number of independent runs limits the precision of the estimated differences between components.

The LogProber method was applied to the Llama prompts to look for signs that the tasks had been memorised during pretraining. All acceleration scores were below the proposed contamination threshold of 1.0. This reduces concern about direct prompt memorisation, but it cannot rule out more general exposure to the tasks or their underlying theories. The highest score was 0.648 for working memory, which was closer to the threshold than the other scores but was not examined further by the authors.

On synthetic datasets with known generating processes, GeCCo recovered the ground-truth model in both the learning and decision-making domains. These tests are valuable because, unlike real behavioural data, the correct mechanism is known. They show that GeCCo can recover a true model under controlled conditions, but they do not guarantee the same result when several mechanisms can explain noisy human behaviour.

Two further checks support the robustness of the pipeline. First, replacing BIC with the Akaike Information Criterion (AIC) produced similar results in the decision-making and planning tasks. Second, replacing the hand-crafted code template with one generated by an LLM produced models with comparable fit. These results reduce concern that the findings depend entirely on one fit statistic or a human-written template.

CENTAUR is a large foundation model trained to predict human behaviour across many cognitive tasks. The authors use it as a high-performance predictive comparison. GeCCo was statistically indistinguishable from CENTAUR in decision making and performed significantly better in learning and working memory. In planning, GeCCo had a numerically worse negative log-likelihood, although the difference was not statistically significant. This comparison shows that small, interpretable models can match a much larger predictor in several tasks; it does not show that GeCCo explains all available behavioural variation.
% ─────────────────────────────────────────────────────────────
\section{Discussion and Critical Evaluation}
The results show that open-source LLMs can serve as proposal engines for computational cognitive modelling. Across four domains, the generated models were competitive with established hand-crafted baselines, and the strongest evidence of improvement came from decision making and working memory. The learning and planning results were more modest. GeCCo also proposed simpler or alternative explanations: a single discount factor linking decision strategies, separate values for experienced and forgone outcomes, and load-dependent reinforcement-learning rates. These proposals are scientifically useful because they suggest hypotheses that researchers may not have considered.

GeCCo may reduce the streetlight effect by expanding the set of models considered, but this claim is only partly tested. The search is still shaped by the prompt, code template, LLM training data, and BIC feedback. It may therefore replace familiar human assumptions with patterns that are common in the LLM's training data. The contamination analysis reduces the likelihood of direct memorisation of the prompts, but it cannot show that the proposed ideas are independent of the existing literature. Human scientific judgement is still required to decide whether a generated mechanism is meaningful.

The use of held-out validation and test data reduces the risk of overfitting, but selection remains an important issue. Each run generates many candidates, and the best candidate is compared with a fixed literature model. The literature model was not developed again with access to the same dataset, time, and computational budget. A fairer test would compare GeCCo with human researchers asked to create new models under matched conditions. The complete discovery process should also be repeated on independent datasets to determine whether the same mechanisms reappear.

The posterior predictive checks provide evidence that the winning models reproduce selected patterns in human behaviour. However, a lower BIC and a good simulation do not establish that a model's internal mechanism is psychologically true. Different mechanisms can produce similar observable choices. This is particularly important for the working-memory result: a load-dependent reinforcement-learning rate and a capacity-limited working-memory process may both account for the same behavioural pattern. A stronger test would design new experiments in which the two accounts make different predictions. The simulation step also used ChatGPT to convert fitting functions into simulation code, adding a further component that should be independently verified.

A practical advantage of GeCCo is that it uses open-source models without fine-tuning. The LLM proposes mathematical structures, while conventional optimisation and held-out BIC determine which candidate is retained. This separation is a sensible safeguard because the LLM does not evaluate its own output. At the same time, the pipeline is not inexpensive in an everyday sense: each run used four A100 GPUs for up to eight hours, and several independent runs were required. Its accessibility is therefore greater than that of some evolutionary searches, but still depends on substantial computing resources.

Template bias is another limitation. The LLM receives a domain-specific function that may restrict the structures it explores, although the follow-up test with an LLM-generated template partly reduces this concern. More importantly, feedback is based mainly on predictive fit. A lower BIC does not automatically identify the best scientific explanation, and a second LLM evaluating theoretical plausibility would not fully solve this problem. Psychological and biological plausibility should be tested through independent experiments rather than judged only through generated text.

Finally, the study covers four controlled laboratory tasks, not the full range of cognitive modelling. It remains unclear whether the pipeline can handle richer data from perception or language, models with many parameters, or natural behaviour outside the laboratory. Cross-task generalisation is particularly important: a useful cognitive theory should explain behaviour across related tasks rather than fit one dataset. Overall, the evidence supports GeCCo as a promising tool for generating candidate models and hypotheses, but not yet as an automated system for establishing cognitive theories.

\subsection*{Context Within Related Work}
Within the paper's literature review, GeCCo is placed alongside a growing set of approaches for automated model discovery in cognitive science. The closest comparison is \citet{castro2025discovering}, who use evolutionary search over symbolic expressions to discover models from human and animal behaviour. Other work translates verbal descriptions of strategies into executable cognitive-model code. GeCCo differs because it uses in-context learning without fine-tuning or a separate evolutionary search algorithm. Its search design is therefore comparatively lightweight and converges within hours rather than days. However, the search is also fairly shallow: ten iterations with three candidates each explore a much smaller space than evolutionary or gradient-based methods that may test thousands of candidates.

More broadly, GeCCo belongs to a family of AI pipelines in which an LLM proposes an output, receives feedback from an external evaluator, and then revises its proposal. GeCCo includes an important safeguard: a candidate is accepted according to BIC or AIC on held-out data, not according to the LLM's own judgement. The refinement loop therefore does not require the LLM to evaluate its own output. The trade-off is that the meaning of ``better'' is largely defined by predictive fit, which does not guarantee that a model is psychologically or biologically meaningful.

\section{Comparison with Long-Term Intracortical BCI Use}

The second paper, by \citet{card2025longterm}, asks a more applied question: can an intracortical brain--computer interface (BCI) provide reliable speech and cursor control during long-term, largely independent use at home? Four microelectrode arrays implanted in the speech motor cortex of one participant with ALS supplied neural signals to a transformer-based brain-to-text decoder and RNN-based cursor decoder. Over 19 months, the participant used the system for more than 3,800 hours and produced 183,060 sentences. In prompted copy tasks, speech decoding exceeded 99\% word accuracy, while during natural use 92.3\% of sentences were rated by the participant as at least mostly correct. These data are especially persuasive for the narrow claim that this particular participant could use the system productively over an unusually long period without daily researcher supervision.

The paper's broader claim of practical clinical viability requires more caution. The study contains only one participant, so it cannot establish that comparable performance will generalise across people, disease stages, implant locations, or patterns of neural-signal degradation. Moreover, the headline accuracy and the real-world evidence measure different things. The 99.2\% figure comes from a constrained copy task, whereas accuracy during spontaneous communication was mainly self-rated and only 53.3\% of utterances were initially decoded completely correctly. The latter measure is ecologically meaningful but subjective, and it is affected by sentence length, fatigue, topic, user corrections, decoder updates, and the language model. The system also remained invasive, used percutaneous wired connections, required trained care partners for approximately 20 minutes of daily setup, and relied partly on gaze control. Thus, the evidence strongly supports long-term feasibility in one case, but not yet scalability or routine clinical independence. A convincing next step would be a preregistered multi-participant study with fixed evaluation periods, objective error metrics for spontaneous speech, comparison against each participant's best existing assistive technology, and measures of fatigue, setup burden, adverse events, and quality of life.

The papers use AI in fundamentally different scientific roles. GeCCo uses an LLM as a \emph{hypothesis generator}: it proposes interpretable mathematical accounts of behaviour, while an external BIC-based evaluator selects among them. Card et al. use deep neural networks as \emph{translational decoders}: the models map high-dimensional neural activity to immediately useful words and cursor movements. GeCCo prioritises interpretability and theory discovery, but its lower BIC does not guarantee that a generated mechanism is psychologically true. The BCI prioritises predictive accuracy and utility, for which a less interpretable transformer is acceptable, but this makes it harder to know which neural features drive errors or whether performance will transfer to a new user. In short, GeCCo has broader task-level coverage but relatively shallow validation within each domain; the BCI has exceptionally deep longitudinal validation but only for one person.

The contrast also clarifies what counts as evidence in each paper. For GeCCo, held-out prediction, recovery of known synthetic models, prompt ablations, and contamination checks support the claim that LLMs can search useful model spaces. They do not by themselves validate the proposed cognitive mechanisms. For Card et al., thousands of hours of voluntary real-world use provide stronger evidence of practical value than a short laboratory benchmark could, but duration within one case cannot substitute for replication across users. Both papers therefore support a bounded claim more strongly than their broadest framing: AI can generate promising candidate theories, and AI decoding can sustain useful communication for one individual. Neither yet establishes general scientific or clinical validity.

\section{Reflection on Reading and Use of AI}

At the beginning of the semester, our reading focused mainly on reconstructing what each model did. In the GeCCo paper, this meant understanding the iterative prompting loop, the separation of training, validation, and test data, and the meaning of BIC and exceedance probability. Preparing the presentation changed our reading strategy: instead of treating every reported improvement as equivalent evidence, we traced each headline claim to the experiment that was supposed to support it. This revealed, for example, that significant BIC improvements were strongest in decision making and working memory, while the evidence in learning and planning was weaker. Reading the BCI paper reinforced the same lesson from the opposite direction. Its 19-month case study provides unusually rich evidence for sustained use by one participant, but the impressive amount of data does not increase the number of independent participants. We therefore learned to separate longitudinal depth from population-level generalisability, and benchmark accuracy from everyday usefulness.

We used generative AI during preparation to clarify technical terminology, check our understanding of statistical metrics and model architectures, translate or rephrase selected passages, and improve the organisation and language of the presentation and report. It also helped us generate candidate criticisms and comparison dimensions. We did not treat these outputs as evidence. Numerical results, methodological details, and the scope of the authors' claims were checked against the original papers, and we rejected or revised suggestions that could not be supported by the text or figures. The final judgments about whether the evidence supports the claims, and what additional experiments would be convincing, are our own.

% ─────────────────────────────────────────────────────────────
\newpage
\bibliographystyle{apalike}
\bibliography{references}

\end{document}
